\chapter{Background}
\label{ch:background}

\section{Partially Observable Decision Making}

A partially observable Markov decision process (POMDP) extends the MDP
formalism with a latent state $s_t$ that is not fully observed
\citep{kaelbling1998planning}. The agent receives observations $o_t$ and must
maintain a belief $b_t(s) = P(s_t = s \mid o_{1:t}, a_{1:t-1})$ over latent
state. Optimal control in POMDPs requires planning over belief space, which is
often intractable in large discrete domains.

Recent work amortizes belief inference with neural sequence models that map
interaction history to action distributions or latent embeddings
\citep{chen2021decision,liu2024tao}. This thesis asks whether
structured auxiliary supervision over interpretable hidden attributes improves
both belief quality and downstream policy robustness.

\section{Simultaneous-Action Games}

In simultaneous-move settings both agents select actions before the environment
advances. The effective decision problem is a partially observed stochastic game.
Unlike turn-based games, the agent cannot observe the opponent's current action
before committing. Policies must therefore encode predictions about opponent
intent from history alone.

\section{Transformers for Sequential Decision Making}

The transformer architecture uses self-attention to relate tokens within a
sequence \citep{vaswani2017attention}. Decision Transformer and related methods apply causal transformers to offline RL by modeling return-conditioned action
sequences \citep{chen2021decision}. In competitive settings, opponent modeling
methods such as TAO embed opponent trajectories for in-context adaptation
\citep{liu2024tao}. The Self-Confirming Transformer generates fictitious belief
tokens representing inferred opponent actions and feeds them back into policy
decoding \citep{zhang2024sct}.

Our approach differs in three ways. We predict structured hidden attributes per
opponent slot rather than only next-opponent-action embeddings. We use a
hierarchical within-turn and cross-turn encoder rather than a flat token
stream. We evaluate combinatorial held-out team and set splits as a primary
metric rather than leaderboard rank alone.

\section{Offline Reinforcement Learning}

Offline RL learns from a fixed dataset without online environment interaction
\citep{sutton2018reinforcement}. Behavioral cloning (BC) imitates expert actions
and provides a stable initialization but suffers from covariate shift.
Conservative Q-Learning (CQL) penalizes overestimation on out-of-distribution
actions \citep{kumar2020conservative}. AMAGO provides another sequence-model
fine-tuning path compatible with transformer backbones.

We use BC with auxiliary belief losses as Stage A and optional CQL or AMAGO
fine-tuning as Stage B. Inference remains a single forward pass with no search.

\section{Combinatorial Generalization}

Humans generalize to novel combinations of known primitives
\citep{lake2018building}. Machine learning systems often fail when test
combinations were absent from training even when individual primitives were
seen. We adopt held-out team keys and held-out move-set tuples as explicit
tests of combinatorial generalization rather than random battle splits alone.

\section{Positioning Relative to Metamon}

\citet{grigsby2025metamon} demonstrated that large sequence models trained on
human Pok\'emon replays can reach strong play without search. That work
establishes feasibility of transformer policies in this domain. Our thesis does
not aim to beat their agents on ladder rank. Instead we study explicit belief
inference, calibration, and OOD degradation under combinatorial splits, which
their paper does not center.
